\subsection{Compact exhaustion with the original finite decay budget}
\label{actual-finite-decay-exhaustion}

The compact kernel estimates require a passage to the prescribed
coefficient with a fixed finite decay exponent.  This passage must
hold at every head, since subsequent arithmetic sums sample individual
points.  Polynomial bounds for the actual quotient coordinates give
that passage with the original dependence on dimension and word length.
The exact cutoff and differential contracts are those in
\cite{finite-decay-source}.

Put
\[
 R=r+1,\qquad d=2r+k+1,\qquad
 X=Nd^2+d^3,\qquad k\ge2.
\]
Retain the original place, normalized tower, top lift, compact type
projections, and all their measures and scalar normalizers.  Assume
the original differential description of the top lift with constant
$C_b$: for each admissible word $w$ and component, its operator on
$G>0$ has order $o$ and polynomial degree $\kappa$, with
\begin{equation}\label{finite-decay-top-budget}
 o\le |C_b|\,|w|d,\qquad
 \kappa\le |C_b|Ro,
\end{equation}
and belongs to the prescribed class \texttt{InDDPlus} with no
internal-minor denominators.  This is precisely the original
\texttt{InDDplus2} hypothesis; its construction is not assumed to
follow from compact exhaustion.

Use the original Frobenius norm
$M(g)=(\sum_{a,b}|g_{ab}|^2)^{1/2}$ and its invariant derivative
seminorms $\|\cdot\|_{A,m}$, as in \JacquetUnfoldingReference.
Constants in estimates below may depend on a fixed place,
head, word, cutoff and representation; their exponent bounds do not.

\begin{theorem}\label{finite-decay-form-limit}
There is a constant $C_c$ depending only on $C_b$ such that the
following holds.  Suppose $f$ is smooth and
$\|f\|_{A,m}<\infty$ for every finite $m$, with the \emph{same}
exponent $A$, and
\[
                       C_cR^4X\le A.
\]
For the original Gram cutoff $F=\chi(G/\epsilon)f$, $\epsilon>0$,
let
\[
 c_n(g)=\operatorname{smoothTransition}
    \left(\frac{4-(n+1)^{-2}\sum_{a,b}|g_{ab}|^2}{3}\right),
       \qquad F_n=c_nF.
\]
Every $F_n$ is an original compact smooth coefficient with a common
positive Gram cutoff.  For every level, every original admissible
word $|w|\le N$, and every pair of heads $a,b$, the prescribed
kernels satisfy
\begin{equation}\label{finite-decay-kernel-limit}
                    \operatorname{RR}_{i,w}(F_n;a,b)
             \longrightarrow \operatorname{RR}_{i,w}(F;a,b).
\end{equation}
The assertion is pointwise at every head, including all relevant
lower dimensional strata.  It supplies the original
\texttt{KB.FormLimit}, with its stated quantifiers and cutoff.
For example, $C_c=10^4(1+|C_b|)^2$ is sufficient.
\end{theorem}

\begin{proof}
We prove quantitative continuity of the actual slices and terminal
orbital vectors, then apply it to this specific exhaustion.

\paragraph{The cutoff and one unit of decay.}
The definition of $\chi$ gives $\delta>0$ such that $\chi(x)=0$
for $x<\delta$, and $\chi(x)=1$ for $x\ge1$.  All its derivatives
are bounded.  Thus all coefficients in question vanish for
$G<\epsilon_0:=\delta\epsilon$.  The Gram polynomial has degree
$2R$.  Invariant differentiation preserves polynomial degree, so
the product rule gives
\begin{equation}\label{finite-decay-gram-cost}
 \|\chi(G/\epsilon)h\|_{B,q}
                       \le C_q\|h\|_{B+2Rq,q}
\end{equation}
for every finite $q$.  The same statement applies to differences.

Every invariant derivative of $c_n$ is bounded uniformly in $n$.
Indeed, the nonconstant derivatives are supported where
$n+1\le\|g\|_{\mathrm F}\le2(n+1)$; each differentiation of
the quadratic argument contributes a factor $(n+1)^{-2}$ and
a polynomial of degree two.  Repeated product and chain rules give
the asserted bound.  On the support of $c_n-1$, or of any of its
derivatives, $\|g\|_{\mathrm F}\ge n+1$.  Hence for every fixed $q$,
\begin{equation}\label{finite-decay-cutoff-convergence}
 \|(c_n-1)f\|_{A-1,q}\le C_q(n+1)^{-1}\|f\|_{A,q}
                                  \longrightarrow0.
\end{equation}
The compact support and common Gram cutoff of $F_n$ follow directly.
No increase of the exponent $A$ with $q$ is used.

\paragraph{Polynomial recovery below the terminal level.}
Fix a head of size $j\le r$, with all block parameters nonzero,
and put $\rho=r-j$.  Use exactly the restricted lower roots,
upper roots and radical coordinates of
Theorem~\ref{partial-head-transfer-actual}.  Write their collection
as $y\in\mathbb R^D$; one may take $D\le2d^2$.  The phase has
absolute value one and is independent of the residual matrix
$\gamma\in G_\rho$.

The inverse formulas \eqref{partial-head-inverse-coordinates} give
radical coordinates of degree at most $2j$ in the entries of the
full matrix $g$, because $\det Z$ is the fixed nonzero head
determinant.  Here and below polynomial means polynomial in real
and imaginary matrix entries; fixed nonzero head determinants
are allowed in its constant coefficients.

For clarity, the block Gaussian step also has a linear, rather than
an iterated exponential, degree bound.  If $Z=L B$, with $L$ block
lower unitriangular and $B=Z_hU$ block upper triangular, let
$I_b$ denote the indices up to the end of block $b$, and $I_{<b}$
the preceding indices.  The exact formulas are
\[
 L_{cb}=\bigl(Z_{c,I_b}Z_{I_b,I_b}^{-1}\bigr)_{b},\qquad
 B_{bc}=Z_{bc}-Z_{b,I_{<b}}Z_{I_{<b},I_{<b}}^{-1}Z_{I_{<b},c}.
\]
All the displayed determinants are products of complete head-block
determinants, hence fixed and nonzero.  Cramer's rule bounds the
degrees of $L,B,U$ by $j+1$; it bounds $L^{-1}$ by $j$ using the
second formula for its strictly lower rows, and $U^{-1}=Z^{-1}LZ_h$
by $2j$.  No internal principal minor inside a head block is used.
Consequently, with a constant depending on this head,
\begin{equation}\label{finite-decay-polynomial-properness}
           1+|y|+M(\gamma)\le C M(g)^{4R}.
\end{equation}
This also proves properness jointly in $y$ and $\gamma$.

A residual left derivative is an ambient left derivative conjugated
by the fixed left radical matrix; a residual right derivative is
conjugated by the right radical matrix.  Each radical matrix and
its inverse has degree at most $2j$ in $g$.  Thus $m$ residual
derivatives cost at most $4jm$ powers of $M(g)$.  The conjugating
matrices are independent of $\gamma$, so no derivatives of them
are inserted in this pullback.  The fixed determinant $D_j$ stays
constant under these residual variations.

\paragraph{The original word operators.}
Expanding the prescribed ordered quantum minors, one letter at
level $j$ has order at most $j$.  Its numerator coefficients have
degree at most $j$, and its only denominators are powers of $D_j$
and $\overline{D_j}$.  An invariant derivative of a polynomial
preserves its degree; differentiating such a denominator adds at
most $j$ to the numerator degree.  Applying the next letter to a
tail therefore increases order by at most $j$ and numerator degree
by at most $j+j^2$.  Induction in the original tail-first order
gives
\begin{equation}\label{finite-decay-word-degree}
        o_w\le |w|j,\qquad
        \kappa_w\le |w|(j+j^2).
\end{equation}
This is an expansion of the actual $\mathsf V_{j,w,e}$, including
the two original signed shifts and both determinant denominators.
The scalar spectral factors and the tensor ordering are unchanged.
For $j=0$, the prescribed zero components vanish and the remaining
components have order zero.

Let $S_{j,w,a,e}h$ be the original normalized partial-head slice
of $\mathsf V_{j,w,e}(\chi(G/\epsilon)h)$.  Combining
\eqref{finite-decay-gram-cost},
\eqref{finite-decay-polynomial-properness} and
\eqref{finite-decay-word-degree}, differentiation under its
\emph{original} integral gives
\begin{equation}\label{finite-decay-slice-bound}
 \|S_{j,w,a,e}h\|_{B,m}\le C\|h\|_{A',m+Nj}
\end{equation}
provided
\begin{equation}\label{finite-decay-slice-budget}
 A'\ge 3NR^2+6Rm+4R(B+2d^2+1).
\end{equation}
Here is the integrability argument in detail.  The differentiated
integrand is bounded by a constant times
\[
 \|h\|_{A',m+Nj}
       M(g)^{-A'+N(j+j^2)+4jm+2R(m+Nj)}.
\]
After multiplication by $M(\gamma)^B$, inequality
\eqref{finite-decay-polynomial-properness} bounds this by
$C\|h\|_{A',m+Nj}(1+|y|)^{-D-1}$ under
\eqref{finite-decay-slice-budget}.  This integrable majorant is
independent of $\gamma$ and of the unit derivative directions.
It proves the asserted seminorm bound, absolute convergence and
all required differentiations.  The exact Gram identity
\eqref{partial-head-gram} gives the common residual cutoff
$\epsilon_0/|\det Z_h|^2$.  The estimate applies to every fixed
head, without restricting its internal minors.

Use the original finite regularity parameters
\[
 B=127R^3X+1,\qquad m=5+3R(N+2)d+Nr.
\]
Since $R\le d$, $d\ge3$ and $X\ge d^3$, we have $m\le8X$.
Thus the right side of \eqref{finite-decay-slice-budget}, plus
one, is at most $600R^4X$.  Our choice of $C_c$ gives this budget
with $A'=A-1$.  Applying \eqref{finite-decay-slice-bound} to
$h_n=(c_n-1)f$ and using
\eqref{finite-decay-cutoff-convergence} proves convergence of
every actual slice in this weighted $C^m$ norm.

\paragraph{Continuity of the original lower-level forms.}
Each entry of the original contracted matrix $tL$ has differential
order at most the residual matrix size and polynomial coefficient
degree at most two: the two vectors depending on $g$ are linear
in $g$ or $g^{-1}$, and $g^{-1}=Jg^*J$ on the unitary group.
Invariant differentiation preserves these degrees.  The exact
residue recursion consequently has order at most $d|w|$ and
coefficient degree at most $2|w|$.  Its old-word rational
corrections are fixed scalars and are retained.

The established Jacquet pairing bound therefore makes the original
form $\mathcal B_w$ continuous on the preceding finite weighted
class.  Explicitly, the threshold for the residual Haar integral
is at most $2(k+\rho)(\rho+1)+1$; our $B$ exceeds it by more
than $2N$, and $m\ge5+dN$.  Apply the pairing bound after the
residue operator, then use the exact sesquilinear identity
\[
 \mathcal B_w(x_n,y_n)-\mathcal B_w(x,y)
   =\mathcal B_w(x_n-x,y_n)+\mathcal B_w(x,y_n-y).
\]
The slice convergence proves \eqref{finite-decay-kernel-limit}
for the original $\operatorname{RR}$ at every $i\le r$.
Any later fixed head weights and word normalizers commute with
this convergence.  No almost-everywhere
representative is substituted at any point.

\paragraph{A quantitative terminal quotient bound.}
For a terminal head $a$ with positive Gram, keep its fixed original
$D$-frame $\gamma_a$.  The exact fibre decomposition in
Theorem~\ref{unitary-density-compact-representative} expresses
its orbital integral using the restricted head roots, the two
pointwise flag kernels $N_0$, and the two original compact groups.
Its noncompact coordinates $y$ again have dimension $D\le2d^2$.
We show
\begin{equation}\label{finite-decay-terminal-properness}
                         1+|y|\le C_a M(g)^{64R^2}.
\end{equation}
This bound concerns the actual quotient, including its discarded
stabilizer roots and its original Haar constants.

First recover the two flag matrices from $Z(g)$ using the preceding
block Cramer formulas with $j=R$.  Their matrices and inverses
are bounded by $C M(g)^{2R}$.  The original Eichler section for
$u=\left(\begin{smallmatrix}A&b\\0&1\end{smallmatrix}\right)$
has blocks
\[
             A,\quad (A^*)^{-1},\quad b,\quad
       b^*(A^*)^{-1}/T,\quad bb^*(A^*)^{-1}/(2T).
\]
The fixed change between $z$ and $z'$ only changes constants.
Since $A$ is unipotent of size $r$, these entries and the inverse
section have degree at most $R+1$ in the flag coordinates.
Stripping both flag actions leaves a matrix $g_0$ with the fixed
reduced matrix $Z_a$, satisfying
\[
                         M(g_0)\le C M(g)^{9R^2}.
\]

Here the positive complement gives a particularly simple bound
for the remaining two $N_0$ factors.  The paired frame
$(f,g_0f')$ has a fixed nondegenerate Gram matrix $H_a$.
Its orthogonal-complement projection is
\[
        \Pi=I-VH_a^{-1}V^*J,\qquad V=(f,g_0f'),
                         \qquad \|\Pi\|\le C M(g_0)^2.
\]
This complement is positive.  If $Q$ is a Euclidean orthonormal
basis of it, then
\[
 \Pi=Q(Q^*JQ)^{-1}Q^*J,\qquad
        \|(Q^*JQ)^{-1}\|=\|\Pi\|.
\]
An isometry from a fixed positive space to this complement is
therefore bounded by $C\sqrt{\|\Pi\|}$.  On the paired-frame
span the corresponding isometry is bounded by $CM(g_0)$ directly.
This bounds every isometry with the prescribed paired-frame action:
the remaining freedom is its isometry on the fixed positive
complement, whose stabilizer is compact.  Thus the bound applies
to every actual fibre representative, not just a chosen factorization.
Apply this to the left pointwise flag stabilizer $\ell$ in
$g_0=\ell\gamma_a u$.  It fixes $f$ and sends $\gamma_af'$
to $g_0f'$, so $\|\ell\|\le CM(g_0)$.  Its inverse has the
same Frobenius bound, since $\ell^{-1}=J\ell^*J$.  The identity
$u=\gamma_a^{-1}\ell^{-1}g_0$ then gives
$\|u\|\le CM(g_0)^2$.  Removing the uniformly bounded original
compact factors gives these same bounds for the two $N_0$
matrices.  Their standard radical coordinates are linear or
quadratic in their matrix entries, hence are bounded by
$CM(g)^{36R^2}$.  This proves the convenient larger bound
\eqref{finite-decay-terminal-properness}.  The measure is exactly
the product root measure and compact Haar measure from the
established fibre decomposition; its fixed positive constants
remain in $C_a$.  No change of orbit or stabilizer measure occurs.

\paragraph{The actual terminal differential class.}
For completeness, the original \texttt{InDDPlus} coefficient is
a finite sum of terms
\[
       P(g)\,D_R(g)^{-a}\overline{D_R(g)}^{-b}G(g)^{-c}
                    \prod_\nu R_{v_\nu}E_\nu(g),
\]
where $\deg P\le\kappa$ and the sum of the derivative lengths
in the fibre factors plus the derivative order applied to the
coefficient is at most $o$.  The prescribed \texttt{FibBound}
gives
\[
        \|R_vE_\nu(g)\|\le C_{v,\epsilon_0}M(g)^{3R|v|}
                         \qquad(G(g)\ge\epsilon_0).
\]
Only $D_R$ and $G$ occur in the denominator, since its set of
internal minors is empty.  They are bounded away from zero on
this domain: $|D_R|^2\ge G\ge\epsilon_0$.  Including the Gram
cutoff derivatives, we obtain
\begin{equation}\label{finite-decay-top-pointwise}
 \|\Theta_{w,e}(\chi(G/\epsilon)h)(g)\|
       \le C\|h\|_{A',o}M(g)^{-A'+\kappa+5Ro}.
\end{equation}
This uses precisely the fibre-derivative allowance in the original
class, rather than treating its differentiated fibre factors as
bounded.

Combine \eqref{finite-decay-top-pointwise} with
\eqref{finite-decay-terminal-properness}.  The original terminal
orbital integral is absolutely convergent and continuous in the
indicated input norm whenever
\[
 A'\ge\kappa+5Ro+64R^2(2d^2+1).
\]
To identify this integral with the original definition for a
noncompact coefficient, use the normalized quotient cut in the
established root-coordinate construction: choose a fixed
integral-one compact bump in the discarded stabilizer roots.
This cut depends on the head and the measures, not on the
coefficient.  The new absolute quotient bound and Tonelli's
theorem show that the cut times the full orbital integrand is
integrable.  Hence the existence branch in the original
$\operatorname{orbJ}$ definition applies, and disintegration gives
exactly this quotient integral with its original normalization.
This extends the compact fibre identity directly; no prior limit
identity is being used.
By \eqref{finite-decay-top-budget}, this right side plus one is
at most $200(1+|C_b|)^2R^4X$.  Thus our choice of $C_c$ gives
the required budget with $A'=A-1$.  Applying the estimate to
$h_n=(c_n-1)f$ proves convergence of each actual terminal vector
$V_{\Theta_{w,e}F_n}(a)$ to $V_{\Theta_{w,e}F}(a)$.
Linearity follows on $G>0$ from the differential description;
the orbital integral stays entirely in that domain.

For a nonempty word, the original terminal kernel is a finite sum
of Hilbert--Schmidt pairings of these vectors after the prescribed
type projections, with the original factors $dH$.  For the empty
word it is the direct Hilbert--Schmidt pairing of $V_F(a)$ and
$V_F(b)$.  The factors are fixed for the heads under consideration,
so the vector convergence proves both kernel limits.  Any subsequent
head weights and word normalization likewise remain unchanged.
Heads of the
wrong size, invalid block heads, and heads with nonpositive Gram
have zero terminal vector by definition.  Invalid lower heads
likewise have zero defining coefficient, and levels $i>r+1$ have
zero original kernel.  This covers every head and level
in the universal quantifier of \eqref{finite-decay-kernel-limit}.
In rank zero the lower quotient has no root variables and the
terminal argument reduces to the original compact quotient; all
the estimates above remain valid.
\end{proof}

The new result removes the finite-decay compact-exhaustion hypothesis
from the local kernel assembly.  It does not construct the arithmetic
top lift or identify a global rational-return sum with the Gram
operator of an integral theta map.  Those retain their full arithmetic
parameters and remain separate obligations.
